% Generated draft: requires mathematical and source review.
\documentclass[11pt]{article}
\usepackage[margin=25mm]{geometry}
\usepackage{fontspec}
\setmainfont{Latin Modern Roman}
\usepackage{amsmath,amssymb,mathtools,unicode-math}
\setmathfont{Latin Modern Math}
\usepackage{enumitem,needspace,xurl,hyperref}
\setlength{\parindent}{0pt}
\setlength{\parskip}{5pt}
\setlength{\emergencystretch}{3em}
\newcommand{\N}{\mathbb{N}}
\newcommand{\Z}{\mathbb{Z}}
\newcommand{\Q}{\mathbb{Q}}
\newcommand{\R}{\mathbb{R}}
\newcommand{\C}{\mathbb{C}}
\newcommand{\F}{\mathbb{F}}
\newcommand{\A}{\mathbb{A}}
\newcommand{\PP}{\mathbb{P}}
\newcommand{\E}{\mathbb{E}}
\newcommand{\eps}{\varepsilon}
\newcommand{\dd}{\,\mathrm{d}}
\begin{document}
\section*{UnsolvedMath ID: 30004740}
\textbf{Source problem code:} catalogue entry
\subsection{Definitions and mathematical statement}
The supplied question asks: \emph{How many \((m,n)\)-rational caustics should a sufficiently smooth strictly convex billiard domain have?}

The following definitions explain a usual planar interpretation; the source record does not explicitly specify the ambient dimension or all conventions needed to make this interpretation definitive. Let \(\Omega\subset\R^2\) be a bounded strictly convex domain with sufficiently smooth boundary. Its billiard trajectories consist of straight segments in \(\Omega\), reflected at \(\partial\Omega\) according to equality of the angles of incidence and reflection. Strict convexity is not here strengthened to everywhere positive boundary curvature.

In this setting, a regular closed curve \(K\subset\Omega\) is called a caustic when reflection of a billiard segment tangent to \(K\) produces another segment tangent to \(K\). The admissible regularity, convexity, and possible degeneracy of \(K\) require specification.

To describe a conventional meaning of the integer pair, orient \(\partial\Omega\), let \(L\) be its perimeter, and lift the successive boundary arclength coordinates of a trajectory to real numbers \(s_j\), choosing \(s_j<s_{j+1}<s_j+L\). A periodic trajectory returning to its initial billiard state after \(n\) reflections has type \((m,n)\) when
\[
 s_n=s_0+mL.
\]
A commonly intended meaning of an \((m,n)\)-rational caustic is a caustic carrying a family of periodic tangent trajectories of this type. The record does not define this term, specify whether every tangent trajectory must be periodic, or distinguish this condition from a condition stated only in terms of rational rotation number. It also does not specify whether \(n\) is the least period or whether \(m\) and \(n\) must be coprime.

Once those conventions and an admissible class of caustics are fixed, the counting quantity can be written
\[
 N_{m,n}(\Omega)
 :=\#\{K:K\text{ is an admissible }(m,n)\text{-rational caustic of }\Omega\},
\]
where curves are counted as geometric subsets and \(\#\) denotes cardinality. This is a possible precise counting formulation, rather than an additional assertion supplied by the record.

\textbf{Unresolved scope.} The source does not say whether \((m,n)\) is fixed, which pairs are allowed, what differentiability threshold is intended, or whether the requested count concerns every domain, a generic domain, or an expectation under a probability distribution. Consequently, it does not determine a unique quantified mathematical assertion. No value, bound, finiteness claim, or genericity claim is asserted here.
\subsection{Short English statement}
The question asks how many rational caustics a sufficiently smooth strictly convex billiard table should possess. Under a usual planar interpretation, these are curves preserved by tangency under billiard reflection and associated with periodic trajectories of prescribed rotation type. The source leaves the precise caustic convention, smoothness threshold, admissible integer pairs, and meaning of the requested count unresolved.
\subsection{Sources}
\begin{itemize}
\item Local UnsolvedMath registry, supplied complete source record: permanent ID \(30004740\), title \emph{Counting Rational Caustics in Convex Billiards}, source problem code \texttt{catalogue entry}, in the collection \emph{Oberwolfach Reports Open Problems}. No individual proposer is named; the supplied proposed year is \(2021\).
\item Supplied report citation: \emph{Dynamische Systeme (2022)}, Oberwolfach Reports, pp.~1735--1803, DOI \texttt{10.4171/owr/2021/33}. Supplied literature-source label: ``Oberwolfach problem report.'' \url{https://doi.org/10.4171/owr/2021/33}. This report reference is supplied as provenance; its contents have not been independently consulted here.
\item Supplied literature-source label: ``density of rational caustics.'' \url{https://arxiv.org/abs/1706.07968}. The record provides no author list, formal title, or precise theorem statement for this reference.
\item Literature assessment dated exactly \texttt{2026-08-22}: this is supplied dated curation, not independently verified. It labels the problem open and states that no general count is known, that fixed-rotation rational caustics can be constructed densely, and that special integrable tables have many caustics. It describes the expected number in general as unresolved. The record explicitly requests independent verification of time-sensitive status and cited claims; no topology or precise density theorem is supplied.
\item The registry labels the statement \texttt{corrected\_verified} and reports a manual semantic, typographic, and LaTeX review. These are supplied editorial assertions. The supplied open label and literature assessment are not independent mathematical verification. The report title uses \(2022\), whereas the proposed year, year tag, and DOI identifier indicate \(2021\); the record does not explain this dating difference.
\end{itemize}
\end{document}
